\begin{definition}
Otoczeniem punktu $x_{0}$ nazywamy przedział otwarty $(x_{0}-\epsilon,x_{0}+\epsilon)$. Przy czym $\epsilon>0$ jest promieniem otoczenia punktu $x_{0}$. Otoczenie punktu $x_{0}$ oznaczamy jako $\mathbb{U}(x_{0},\epsilon)$.
\end{definition}
\begin{remark}
Zauważmy, że punkt $x_{0}$ należy do swojego otoczenia.$\\$
Zauważmy też, że $x\in\mathbb{U}(x_{0},\epsilon)\iff x_{0}-\epsilon< x< x_{0}+\epsilon$. Czyli $x-x_{0}>-\epsilon\wedge x-x_{0}<\epsilon$ co z własności wartości bezwzględnej daje $|x-x_{0}|<\epsilon$.$\\$
W szczególności $x_{0}\in\mathbb{U}(x_{0},\epsilon)\iff x_{0}> x_{0}-\epsilon\wedge x_{0}< x_{0}+\epsilon$. Zatem $x_{0}-\epsilon< x_{0}< x_{0}+\epsilon\iff -\epsilon< x_{0}-x_{0}<\epsilon$. Czyli $-\epsilon<0<\epsilon$. Wszystko się zgadza, ponieważ $\epsilon> 0$.    
\end{remark}
\begin{definition}
Ciągiem liczbowym nazywamy funkcję, która odwzorowuje dowolny podzbiór zbioru liczb naturalnych w zbiór liczb rzeczywistych.    
\end{definition}
\begin{remark}
Jeżeli dziedziną ciągu jest podzbiór niewłaściwy zbioru liczb naturalnych to ciąg ten jest ciągiem liczbowym nieskończonym. Podobnie, gdy dziedziną ciągu jest dowolny nieskończony podzbiór właściwy zbioru liczb naturalnych. W przypadku zaś, gdy dziedziną ciągu jest dowolny skończony podzbiór właściwy zbioru liczb naturalnych to ciąg ten jest ciągiem liczbowym skończonym.$\\$
Należy podkreślić, że liczby $0$ nie włączam do zbioru liczb naturalnych. Zatem przyjmuję, że $\mathbb{N}=\{1,2,3,...\}$.$\\$
Można zatem zapisać, że ciąg jest funkcją $f:\mathbb{D}\rightarrow\mathbb{R}$, gdzie $\mathbb{D}\subseteq\mathbb{N}$.$\\$
Wówczas funkcję $f$ oznaczamy jako $(a_{n})$, przy czym $n\in\mathbb{D}$.
\end{remark}
\begin{definition}
Wartość funkcji $f:\mathbb{D}\rightarrow\mathbb{R}$ dla argumentu $n\in\mathbb{D}\subseteq\mathbb{N}$ nazywamy $n$-tym wyrazem ciągu $(a_{n})$, co zapisujemy $f(n)=a_{n}$.
\end{definition}
\begin{remark}
Należy podkreślić, że w przypadku ciągu ważna jest kolejność jego wyrazów. Wówczas każdy ciąg skończony ma wyraz pierwszy i ostatni, zaś ciąg nieskończony ma wyraz pierwszy ale nie ma wyrazu ostatniego.    
\end{remark}